\documentclass[11pt]{beamer}
\usetheme{Madrid}
\usefonttheme{serif}

\usepackage[utf8]{inputenc}
\usepackage[spanish]{babel}
\usepackage[T1]{fontenc}

\usepackage{amsmath}
\usepackage{amsfonts}
\usepackage{amssymb}
\usepackage{graphicx}
\usepackage{listings}

\DeclareMathOperator{\sen}{sen}
\DeclareMathOperator{\tg}{tg}

\setbeamertemplate{caption}[numbered]

%\setbeameroption{show notes on second screen=right}  % También puedes usar: show notes on second screen=right

\author[LLO]{Ing. Leonardo Luis Ortiz}
\title{Introducción a la programación RTL en C++ para simulación de procesamientos digitales de señales}
% Informe su email de contato o comando a seguir
% Por ejemplo, alcebiades.col@ufes.br
\newcommand{\email}{leonardo.ortiz@ib.edu.ar}
%\setbeamercovered{transparent} 
\setbeamertemplate{navigation symbols}{} 
\logo{\includegraphics[scale=0.08]{images/IB.png}} 
\institute[]{Instituto Balseiro \par Centro Atómico Bariloche} 
\date{\today} 
%\subject{}

% ---------------------------------------------------------
% Selecione un estilo de referencia
\bibliographystyle{apalike}

%\bibliographystyle{abbrv}
%\setbeamertemplate{bibliography item}{\insertbiblabel}
% ---------------------------------------------------------

% ---------------------------------------------------------
% Incluir los slides 
%\usepackage[spanish,hyperpageref]{backref}

%\renewcommand{\backrefpagesname}{Cita de página(s):~}
%\renewcommand{\backref}{}
%\renewcommand*{\backrefalt}[4]{
	%	\ifcase #1 %
	%		Sin citas en el texto.%
	%	\or
	%		Citado en la página #2.%
	%	\else
	%		Citado #1 veces en las páginas #2.%
	%	\fi}%
% ---------------------------------------------------------
%% Configuración para el código
%\lstset{
%	language=C++,
%	basicstyle=\ttfamily\scriptsize,
%	keywordstyle=\color{blue},
%	commentstyle=\color{gray},
%	stringstyle=\color{orange},
%	showstringspaces=false,
%	breaklines=true,
%	frame=single,
%	numbers=left,
%	numberstyle=\tiny,
%	tabsize=2
%}


\begin{document}
	
	\begin{frame}
		\titlepage
	\end{frame}
	
	\begin{frame}{Agenda}
		\tableofcontents 
	\end{frame}
	
	\section{Introducción}
	\begin{frame}
		\begin{itemize}
			\item Hardware Description Language - HDL 
			\item Event-Driven Programming - EDP
		\end{itemize}
	\end{frame}
	
	\section{HALCON simulator}
	\begin{frame}
	\end{frame}
	
	\section{Example 1. Signal \& Noise Generator. Filter}
	\begin{frame}
	\end{frame}
	\section{Example 2. Frequency Domain Analysis. FFT/IFFT}
	\begin{frame}
	\end{frame}
	\section{Example 3. fftw library. Update example 2}
	\begin{frame}
	\end{frame}
	\section{Example 4. Audio processing or some type of modulation}
	\begin{frame}
	\end{frame}
	\section{Example 5. Signal Pattern Generator. Characteristics}
	\begin{frame}
	\end{frame}
	\section{Final project day 1}
	\begin{frame}
	\end{frame}
	\section{Final project day 2}
	\begin{frame}
	\end{frame}
	
	\section{Final project day 3}
	\begin{frame}
	\end{frame}



\section{Introducción}

\begin{frame}{¿Por qué procesamiento de señales?}
	La siguientes áreas de Ciencia e Ingeniería	se han beneficiado por el rápido crecimiento y avances de las técnicas de procesamiento de señales. 
	\vspace{0.2cm}
	\begin{itemize}
		\item Machine Learning
		\item Análisis de Datos
		\item Computer Vision
		\item Procesamiento de Imágenes e Imágenes Médicas
		\item Systemas de Comunicación
		\item Electrónica de Potencia
		\item Problabilidad y Estadística
		\item Análisis Numérico
		\item Teoría de Decisión
		\item Circuitos Integrados
	\end{itemize}
\end{frame}

\begin{frame}{Procesamiento Digital de Señales - DSP}
	Contenidos mínimos necesarios para DSP	\newline
	\begin{itemize}
		\item Fundamentos de procesamiento de señales
		\item Conversión Analógica-Digital
		\item Sampleo y Reconstrucción
		\item Teorema de Nyquist
		\item Convolución
		\item Eliminación de ruido en señales
		\item Transformada de Fourier
		\item Diseño de filtros FIR e IIR 		
	\end{itemize}
\end{frame}

\begin{frame}{Herramientas de Simulación}
	\begin{tikzpicture}[remember picture,overlay]
		\node at (2.5,1.8) 
		{\includegraphics[width=0.5\textwidth]{images/Simulation_tools/Ansys_Lumerical.png}};
		\node at (2.5,-2.5) {\includegraphics[width=0.4\textwidth]{images/Simulation_tools/siemens-electric-motors-design-and-simulation.png}};
		\node at (9,2) {\includegraphics[width=0.55\textwidth]{images/Simulation_tools/Ansys.png}};
		\node at (9,-2.5) {\includegraphics[width=0.5\textwidth]{images/Simulation_tools/PSpice_Product.png}};
		
	\end{tikzpicture}	
	
\end{frame}

\begin{frame}{Lenguajes de Descripción de Hardware - HDL}
	En la actualidad se vuelve indispensable diseñar, modelar y simular el funcionamiento del hardware, principalmente en la industria ya que de este modo se evitan grandes pérdidas por errores de diseño. Para poder lograr esto existen los "Lenguajes de Descripción de Hardware".\newline
	
	Los Lenguajes de Descripción de Hardware (HDL) son una metodología de diseño digital de hardware a través de software.
	
	\note{
		HDL incluye explícitamente la noción de tiempo.
	}
\end{frame}

%	\begin{frame}[fragile]{Diseño digital}
	%	\begin{tikzpicture}[node distance=2cm]
		%		
		%		% Módulo 1
		%		\node[module, minimum width=5cm, minimum height=3cm] (m1) {};
		%		\node[cloud] (c1) at (-1,0) {Cloud 1};
		%		\node[block, right=1cm of c1] (b1) {};
		%		\node[below=1.2cm of c1] {Module 1};
		%		
		%		% Módulo 2
		%		\node[module, right=0.2cm of m1, minimum width=5cm, minimum height=3cm, anchor=west] (m2) {};
		%		\node[cloud, right=1cm of b1] (c2) {Cloud 2};
		%		\node[block, right=1cm of c2] (b2) {};
		%		\node[below=1.2cm of c2] {Module 2};
		%		
		%		% Módulo 3
		%		\node[module, below=0.2cm of m1.south west, anchor=north west, minimum width=10.2cm, minimum height=3.5cm] (m3) {};
		%		\node[cloud, below=1.5cm of c2] (c3) {Cloud 3};
		%		\node[block, below=1.5cm of c3, minimum width=3cm] (b3) {};
		%		\node[below=1.8cm of c3] {Module 3};
		%		
		%		% Conexiones
		%		\path[line] (c1.east) -- (b1.west);
		%		\path[line] (b1.east) -- (c2.west);
		%		\path[line] (c2.east) -- (b2.west);
		%		
		%		\path[line] (b1.south) |- (c3.west);
		%		\path[line] (b2.south) |- (c3.east);
		%		
		%		\path[line] (c3.south) -- (b3.north);
		%		
		%		% Entradas y salidas
		%		\path[line] ([xshift=-1cm]c1.west) -- (c1.west);
		%		\path[line] (b2.east) -- ++(1cm,0);
		%		\path[line] (b2.east) -- ++(1cm,-0.5cm);
		%		
		%	\end{tikzpicture}
	%	\end{frame}

\begin{frame}{Diseño digital}
	\begin{figure}[htb]
		\centering
		\includegraphics[width=0.9\textwidth]{images/Simulation_tools/pspice_cadence.png}
	\end{figure}
\end{frame}

\begin{frame}{Diseño digital}
	\begin{figure}[htb]
		\centering
		\includegraphics[width=0.5\textwidth]{images/System_level_design_components}
	\end{figure}
\end{frame}

\section{Lenguajes de Descripción de Hardware - HDL}

\begin{frame}{Lenguajes de Descripción de Hardware - HDL}
	En la actualidad se vuelve indispensable diseñar, modelar y simular el funcionamiento del hardware, principalmente en la industria ya que de este modo se evitan grandes pérdidas por errores de diseño. Para poder lograr esto existen los "Lenguajes de Descripción de Hardware".\newline
	
	Los Lenguajes de Descripción de Hardware (HDL) son una metodología de diseño digital de hardware a través de software.
	
	\note{
		HDL incluye explícitamente la noción de tiempo.
	}
\end{frame}

\begin{frame}{Lenguajes de Descripción de Hardware - HDL}
	\textbf{Ventajas} \newline
	
	\begin{itemize}
		\item INDEPENDIENTE DE LA TECNOLOGÍA\newline
		\item SOPORTA VARIOS ESTILOS DE DESCRIPCIÓN (Behavioral, Dataflow, Structured, Switch)\newline
		\item AHORRO DE TIEMPO Y COSTOS\newline
	\end{itemize}
	
	
	\note{
		INDEPENDIENTE DE LA TECNOLOGÍA: el mismo modelo puede ser sintetizado en librerías de distintos fabricantes.
		SOPORTA 3 ESTILOS DE DESCRIPCIÓN: (Behavioral, Dataflow, Structured)
		AHORRO DE TIEMPO Y COSTOS: facilita las correccionesen el diseño debidas a fallos de diseño o cambios en las especificaciones.
		
	}
\end{frame}

\begin{frame}{Lenguajes de Descripción de Hardware - HDL}
	\textbf{Módulos} 
	
	El módulo es la unidad básica de diseño, en la que se describe el funcionamiento de alguna unidad lógica. Puede existir un sólo módulo o un conjunto de módulos que se instancian con un módulo estructural.\newline
	
	
	\textbf{Entradas y Salidas} 
	
	Un módulo define la información que describe la relación entre las entradas y salidas de un circuito lógico. Las entradas y salidas están definidas por las palabras reservadas según sea el caso del lenguaje utilizado.\newline
	
	\textbf{Primitivas} 
	
	Elementos lógicos predefinidos, que son elementos estructurales que se pueden instanciar en un diseño más grande para formar una estructura más compleja.(Ej. en Verilog: and, or, xor, not)
	
	
	\note{
		Un mpodulo Verilog es una representación de software de la estructura y comportamiento del hardware físico.	
	}
	
\end{frame}

	
\end{document}